\section{Recursive forcing: construction and transfer proofs}
\label{sec:annex-p06-forcing-recursion-truth}

\begin{annexitem}{annex:p06.ftr.name-foundations}{thm:p06-ftr-name-foundations}{Set recursion before forcing}
\begin{proof}
\paragraph{Physical ranks without an assumed name recursion.}
For any set \(x\), form \(T=\operatorname{tc}(\{x\})\) by the sequence
\(T_0=\{x\}\), \(T_{n+1}=T_n\cup\bigcup T_n\), followed by its union.
The completed ordinal recursion theorem
\cref{thm:p06-definable-ordinal-recursion}, restricted to \(\omega\),
and Replacement and Union construct this transitive set.
Put \(T_y=\operatorname{tc}(\{y\})\) for \(y\in T\).
A rank approximation on a transitive subset \(U\subseteq T\) is an
ordinal-valued function \(r\) satisfying
\[
 r(y)=\sup\{r(z)+1:z\in y\}\qquad(y\in U).
\]
Two approximations agree on their common domain. Otherwise Separation forms
the set of disagreement points, and Foundation supplies a disagreement
\(y\) with no member at which they disagree. Both domains are transitive,
so both contain every \(z\in y\); their two suprema are identical.

Foundation induction proves that an approximation exists on every \(T_y\).
If approximations exist on all \(T_z\), \(z\in y\), uniqueness makes
the formula assigning that graph to \(z\) functional. Replacement on the
set \(y\) and Union give their compatible union \(g\). Its domain is
\(U=\bigcup_{z\in y}T_z\), and \(T_y=\{y\}\cup U\).
If \(y\in U\), this graph already restricts to the required approximation.
Otherwise append \(\langle y,\sup_{z\in y}(g(z)+1)\rangle\).
The new value is an ordinal by Replacement on \(y\) and Union; old
equations are unchanged because \(U\) is transitive. This proves the
induction step. Formally, Foundation on the set of failures in \(T\)
would give a failure all of whose members satisfy that step.
Replacement on \(T\) now collects these unique approximations, and their
union is a rank function on \(T\). Uniqueness also gives agreement for
different bounding transitive closures. In particular membership strictly
decreases physical rank. Foundation induction using the defining clauses
of \(V\), whose existence is proved in
\cref{thm:p06-hierarchy-recursion-existence}, gives
\(y\in V_{r(y)+1}\). Thus every set has a rank bound; this argument has
not assumed the earlier sketched hierarchy theorem.

\paragraph{A bounded subname closure.}
From each ordered pair in a set \(\tau\) take its unique first coordinate,
discarding elements that are not ordered pairs. Pair identity makes this
a definable partial projection; Separation on \(\tau\) and Replacement
on its accepted elements form its image. Starting at \(\{\sigma\}\),
adjoin these images at successive finite stages. Each subname lies in
\(\operatorname{tc}(\{\sigma\})\), since the Kuratowski pair containing
it gives a finite membership path. The union of the stages is therefore
a set \(S_\sigma\). Testing every entry of every member of this set
for being a pair with second coordinate in \(P\) is a single first-order
formula. Malformed entries make this test false; they are not new names.
For an accepted entry \(\langle\pi,p\rangle\in\tau\), its physical
rank calculation gives \(r(\pi)<r(\tau)\).

On any subname-closed set of accepted names, construct a rank table by
ordinal recursion on their bounded physical ranks. At a successor, assign
to every name of the newly reached physical rank the supremum of the
already assigned child-name ranks plus one; at a limit take the union of
the previous compatible tables. The names assigned at one step form a
subset of the bounding closure; Replacement gives their ordinal outputs
and graph. Lower physical ranks contain all children, so each assignment
has all of its inputs. The recursion operation is total on arbitrary
history tables by taking zero for a missing lookup; the completed ordinal
recursion theorem supplies the table, and induction shows that on this
actual construction the missing-lookup branch is never used.
Least physical rank of a disagreement proves uniqueness. The same
induction proves compatibility between bounding closures and
\(\rho_P(\sigma)\leq r(\sigma)\).

The \(W\)-hierarchy is constructed by that same completed ordinal
recursion theorem, with successor-operation value
\(\mathcal P((\bigcup\operatorname{ran}(f))\times P)\) when the
history has ordinal domain. Cartesian products are bounded by
\(\mathcal P(\mathcal P(a\cup b))\) before Separation; no class-sized
product is formed. This is a unique set value for every history, including
domain zero. The stages are increasing: their earlier-name unions increase,
and inclusion of the product bases implies inclusion of the corresponding
power sets. Induction on the intrinsic name rank gives
\(\sigma\in W_{\rho_P(\sigma)}^P\): every child is in its own earlier
stage. Conversely, if \(\sigma\in W_\alpha^P\), every child is in
some \(W_\beta^P\), \(\beta<\alpha\); induction gives its rank at
most \(\beta\), and hence \(\rho_P(\sigma)\leq\alpha\).
This also proves agreement of the closure and hierarchy definitions.

Each construction is first-order definable, with set parameters, and can
be carried out inside \(\mathcal M\). Its computed closure, pair tests,
ordinal suprema and physical-rank equations are correct externally by
transitivity. The uniqueness arguments show that the ground rank and
closure agree with the ambient ones. This is not an assertion that ground
power sets contain every ambient subset.

\paragraph{Evaluation and canonical names.}
With \(G\) as an ambient set parameter, the same bounded physical-rank
tables assign to each name the set of values of children whose tags belong
to \(G\). Replacement on the set of children and Separation on its
tagged entries form each value. The ordinal table construction supplies
existence; least physical-rank disagreement supplies uniqueness and
compatibility. Thus evaluation is available without forcing or genericity.
For \(x\in M\), the analogous ground recursion assigning
\(\{\langle\check y,1_P\rangle:y\in x\}\) gives \(\check x\in M\).
Membership induction proves \((\check x)^G=x\) when \(1_P\in G\).
Replacement inside \(M\) forms
\(\dot G=\{\langle\check p,p\rangle:p\in P\}\), whose value is \(G\).

The ground names are an ambient subset of the set \(M\), and ambient
Replacement on this set with the unique evaluation formula forms
\(M[G]\). If \(a\in\sigma^G\), its witnessing subname belongs to \(M\)
by transitivity, so \(a\in M[G]\); hence \(M[G]\) is transitive.
Induction on name rank proves
\[
 \operatorname{rank}(\sigma^G)
 =\sup_{a\in\sigma^G}(\operatorname{rank}(a)+1)
 \leq\sup_{\pi\in A_\sigma}(\rho_P(\pi)+1)=\rho_P(\sigma).
\]
The already constructed physical rank of an ordinal \(\alpha\) is
\(\alpha\), by induction and the supremum of its predecessors plus one.
If an ordinal \(\alpha=\sigma^G\) belongs to \(M[G]\), then
\(\alpha\leq\rho_P(\sigma)\in M\). Equality or transitivity therefore
puts \(\alpha\) in \(M\). The converse follows from canonical names.
The bounded ordinal predicate of \cref{thm:p06-ordinal-foundation-facts}
is absolute for these transitive sets, so this is also the ordinal
comparison used after \(\mathsf{ZFC}\) is established in the extension.

Karagila's continuous hierarchy has empty stage zero and
condition--name pairs. Recursively reverse each pair and transform its
subname. The same bounded table argument supplies this map and its inverse;
induction gives preservation of intrinsic rank and evaluated value.
His stage \(\alpha+1\), consisting of names of rank at most \(\alpha\),
corresponds to this \(W_\alpha^P\). At a source limit stage the ranks
are strictly below the limit. The reindexing and pair reversal are
explicit; no literal equality of the differently written stage systems
or of the paid source versions has been assumed.
\end{proof}
\end{annexitem}

\begin{annexitem}{annex:p06.ftr.recursion}{thm:p06-ftr-recursion}{The actual bounded atomic tables and formula definitions}
\begin{proof}
Fix an ordinal \(\alpha\in M\) and work in \(\mathcal M\) on the
set \(W_\alpha^P\). It is closed under subnames. All its names have
intrinsic rank at most \(\alpha\). The sorted-rank pairs
\((u,v)\), \(v\leq u\leq\alpha\), have a well-founded lexicographic
order: any nonempty subset has a least first coordinate and then a least
second coordinate. If one coordinate of an unsorted pair is lowered,
the sorted pair either has smaller maximum, or the same maximum and
smaller minimum. Every recursive call in the atomic definitions therefore
refers to an earlier sorted-rank pair.

Construct a table recording two subsets of \(P\) for every ordered
name pair \((\sigma,\tau)\in(W_\alpha^P)^2\), one for equality and
one for membership. There is an explicit nested ordinal recursion:
the outer recursion visits \(u\leq\alpha\), and the inner recursion
at that \(u\) visits \(v\leq u\). At \((u,v)\), form by Separation
the truth bits for all name pairs of that sorted rank from the bounding
set \(P\times(W_\alpha^P)^2\times\{0,1\}\).
The equality bit tests both displayed \(\mathsf{Incl}\) clauses, looking
up the lower membership bits for each subname and strengthening.
The membership bit tests density of the displayed witness set, looking
up lower equality bits. Density expands to
\(\forall q\leq_P p\,\exists t\leq_P q\) with the witness condition
on \(t\); hence these are finite first-order specifications with
bounded poset and entry quantifiers. Missing history lookups have the
fixed default false, making each auxiliary operation total even on an
arbitrary history. On the actual table, every needed lookup is earlier
and already present. Two applications of
\cref{thm:p06-definable-ordinal-recursion}, and Replacement collecting
the inner tables, construct the full set table. At the zero pair,
membership in the empty name is false and both inclusion tests for
equality of empty names are vacuous; this fixes the base explicitly.

Induction through the nested order proves uniqueness of every bit, since
the right sides use only already unique earlier bits. The same induction
on a pair in two bounds \(\alpha,\beta\) proves compatibility: all
subnames and their already compared bits agree. For any two ground names,
take a bound at least their ranks. A first-order formula saying that there
are this bound and the specified nested-recursion tables with the indicated
bit therefore defines the atomic relation. Compatibility prevents a
different bound from changing its answer. This proves both existence and
definability, not merely a name for an unperformed recursion.

For a fixed standard formula, proceed externally through its finite
syntax. Falsehood has the always-false defining formula. Negation replaces
the previous formula by its universal no-strengthening clause. Conjunction
conjoins the two previous formulas. Existential quantification applies
the density formula to \(\exists\sigma\,(\mathsf{Name}_P(\sigma)\land
\mathsf{Frc}_{P,\theta}(q,\sigma,\bar\tau))\).
The names quantified here are ordinary sets in the ground structure;
Separation on \(P\) produces the condition set despite this unbounded
set quantifier. Fresh-variable renaming at each insertion prevents
capture. The three displayed Boolean/quantifier expansions are finite
syntactic operations and have their standard classical satisfaction
equivalences. Thus all of the active primitives receive definitions.
Uniqueness follows from the atomic uniqueness and these fixed clauses.

This construction supplies one ground defining formula for each fixed
standard formula, or for a fixed finite family. It has not quantified
internally over all formula codes. In particular the one-condition poset
would not turn the construction into a truth definition for the universe.
Neither generic filters nor \(M[G]\) occurred in the construction.
\end{proof}
\end{annexitem}

\begin{annexitem}{annex:p06.ftr.density}{thm:p06-ftr-density}{All monotonicity and density cases}
\begin{proof}
We suppress the fixed subscripts on recursive forcing. First prove atomic
monotonicity and density closure together by induction on sorted name ranks.
For membership, restriction of its dense witness set from the cone below
\(p\) to the cone below a strengthening proves monotonicity.
For equality, every inclusion test at a strengthening is one of the
tests already imposed at \(p\). This proves monotonicity there too.

Suppose conditions forcing \(\sigma\in\tau\) are dense below \(p\).
Given \(r\leq_P p\), first strengthen to a condition forcing that
membership, and then strengthen into its defining witness set. The
result is a defining membership witness below \(r\). Thus the witness
set is dense below \(p\), proving density closure for membership.
For equality, fix \(\langle\pi,s\rangle\in\sigma\) and
\(r\leq_P p,s\). Below every \(t\leq_P r\), the assumed density gives
\(q\leq_P t\) forcing equality. Its inclusion test gives
\(q\forces\pi\in\tau\). Conditions forcing this lower-rank membership
are therefore dense below \(r\); its already proved density closure
gives \(r\forces\pi\in\tau\). Do this for every entry, then do the
symmetric argument for \(\tau\). Both inclusion tests hold at \(p\).
Conversely, monotonicity makes the whole cone below any forcing condition
consist of forcing conditions, which proves the other density direction.
Rank induction also gives \(p\forces\sigma=\sigma\) for every \(p\):
an entry \(\langle\pi,s\rangle\) witnesses \(\pi\in\sigma\) throughout
the cone below \(p,s\), using lower-rank \(\pi=\pi\).
Symmetry is immediate from the two inclusion tests.

Now induct on primitive formula formation. Falsehood is never forced;
on a nonempty cone its forcing set cannot be dense. For conjunction,
monotonicity passes to both conjuncts. If conjunction conditions are dense,
each conjunct's forcing conditions are dense, so both hold at \(p\) by
the induction hypothesis. For negation, shrinking the cone preserves the
no-positive-strengthening clause. If negation conditions are dense below
\(p\) but some \(r\leq_P p\) forces the positive formula, a negation
condition below \(r\) also forces the positive formula by monotonicity,
contradicting its no-positive-strengthening clause. Thus \(p\) forces
the negation. A condition cannot force both signs by that same clause.

For an existential, restricting its dense witness set proves monotonicity.
If existential-forcing conditions are dense below \(p\), below any given
\(r\leq_P p\) first find such a condition, then enter its witness set.
The resulting condition forces the matrix at a particular ground name.
Hence the original existential witness set is dense below \(p\).
This proves density closure without choosing one name for all conditions.

For any already defined formula, its decision set is dense: below a
condition either a strengthening forces it, or that condition itself
forces its negation. Monotonicity makes this set open. The fixed defining
formulas and Separation on \(P\) place it in \(M\). In particular,
if \(p\) does not force a formula, density closure gives
\(q\leq_P p\) below which no condition forces it; this \(q\) forces
its negation. This conclusion is combinatorial, not semantic.

For disjunction, if the positive-branch conditions are dense below \(p\),
no strengthening can force both negative branches. Otherwise a further
positive-branch condition contradicts one of them. Thus
\(p\forces\neg(\neg\varphi\land\neg\psi)\).
Conversely, if that branch set is not dense, there is \(q\leq_P p\)
with no positive-branch strengthening. Then \(q\) forces both negations,
contradicting the supposed expanded disjunction at \(p\).

For implication, if \(p\forces\neg(\varphi\land\neg\psi)\) and
\(q\leq_P p\) forces \(\varphi\) but not \(\psi\), the preceding
no-forcing conclusion gives \(r\leq_P q\) forcing \(\neg\psi\).
It still forces \(\varphi\), a contradiction. Conversely, if every
\(\varphi\)-forcing strengthening forces \(\psi\), none can force
\(\varphi\land\neg\psi\). This is exactly the expanded implication.

Finally \(p\forces\neg\exists x\neg\theta(x)\) holds iff no
\(q\leq_P p\) forces a negated instance: such an instance would make
its existential forced on its whole cone, and any forced existential
has such a witness-strengthening. For each fixed ground name, absence
of a negated-instance strengthening makes the positive-instance conditions
dense by decision; density closure forces that instance at \(p\).
Conversely, all instances at \(p\) exclude every negated-instance
strengthening. This proves the universal clause. Monotonicity and
density for the expanded constructors follow either from these clauses
or from their just-proved primitive expansions. No constructor is omitted.
\end{proof}
\end{annexitem}

\begin{annexitem}{annex:p06.ftr.atomic-truth}{thm:p06-ftr-atomic-truth}{Membership, equality and the disagreement dense set}
\begin{proof}
\paragraph{Meeting a local dense set.}
If \(D\in M\) is dense below \(p\in G\), then
\[
 D\cup\{q\in P:q\text{ is incompatible with }p\}
\]
is globally dense and belongs to \(M\). A condition compatible with
\(p\) first strengthens to their common extension, then into \(D\);
an incompatible condition already belongs to the second part.
Genericity meets this set. Directedness and \(p\in G\) rule out the
incompatible part, so \(G\cap D\ne\varnothing\).
All subsequent witness and decision sets are in \(M\) by the preceding
definability theorem and Separation; this local meeting argument applies
without any external enumeration of \(M\).

Prove the two positive atomic equivalences and their negations together
by induction on the sorted-rank pair. Suppose
\(p\in G\) forces \(\sigma\in\tau\). Its defining witness set is
dense below \(p\). A member \(q\in G\) of that set has
\(\langle\pi,r\rangle\in\tau\), \(q\leq_P r\), and
\(q\forces\sigma=\pi\). The equality pair has lower measure, so
\(\sigma^G=\pi^G\). Upward closure puts \(r\) in \(G\), and
evaluation gives \(\sigma^G\in\tau^G\).
Conversely, membership of evaluated names gives such an entry with
\(r\in G\) and \(\sigma^G=\pi^G\). The lower equality induction
gives \(q\in G\) forcing that equality. Directedness gives
\(t\in G\), \(t\leq_P q,r\). Every strengthening of \(t\) is a
membership witness, so \(t\forces\sigma\in\tau\).

Suppose now \(p\in G\) forces \(\sigma=\tau\). For
\(a\in\sigma^G\), take an entry \(\langle\pi,r\rangle\in\sigma\)
with \(r\in G\) and \(a=\pi^G\). A common strengthening
\(q\in G\) of \(p,r\) passes the inclusion test, so
\(q\forces\pi\in\tau\). This is lower-rank membership and gives
\(a\in\tau^G\). Interchanging the names proves equality of their values.

\paragraph{The converse equality direction.}
Define \(B_{\sigma,\tau}\) to consist of conditions \(q\) for which
either there is \(\langle\pi,r\rangle\in\sigma\) with
\(q\leq_P r\) and \(q\forces\neg(\pi\in\tau)\), or there is the
symmetric witness from \(\tau\). Then
\[
 D_{\sigma,\tau}
 =\{q:q\forces\sigma=\tau\}\cup B_{\sigma,\tau}
\]
is dense. Indeed, if \(p\) does not force equality, one inclusion
test fails. It gives an entry \(\langle\pi,r\rangle\) and
\(s\leq_P p,r\) not forcing the relevant lower membership.
Lower-rank density closure gives \(q\leq_P s\) with no positive
membership strengthening, hence forcing its negation. Thus
\(q\in B_{\sigma,\tau}\). The other inclusion failure is symmetric.
This defines one ground dense set; it is not a purported infinite
intersection of conditions separately found in \(G\).

If \(\sigma^G=\tau^G\), genericity meets \(D_{\sigma,\tau}\).
A meeting condition in \(B_{\sigma,\tau}\) would put its tag \(r\)
in \(G\), making \(\pi^G\) an element of one evaluated name.
The lower negated-membership induction makes it absent from the other,
contradicting equality. Therefore the meeting condition forces equality.
This proves the previously unresolved positive equality direction.

At the present rank, a condition in \(G\) forcing the negation of an
atom makes the atom false: if it were true, the positive equivalence
just proved would give another condition in \(G\) forcing it, and a
common strengthening contradicts the negation clause. Conversely, if
the atom is false, the ground decision dense set is met by \(G\);
its positive alternative is impossible by soundness, leaving the
negation alternative. This completes the joint induction, including
empty names and zero ranks.

\paragraph{Canonical ground atoms used in the boundary example.}
Induction on the same rank measure gives, for \(x,y\in M\) and every
\(p\in P\),
\[
 p\forces\check x\in\check y\ \Longleftrightarrow\ x\in y,
 \qquad p\forces\check x=\check y\ \Longleftrightarrow\ x=y.
\]
For membership, entries of \(\check y\) have form
\(\langle\check z,1_P\rangle\), \(z\in y\). If \(x\in y\), the
entry at \(z=x\), with lower reflexive equality, witnesses density
throughout the cone. If \(x\notin y\), lower equality is false at
every such entry, so the witness set is empty. For equality of equal
ground sets, all inclusion tests are these true lower membership
instances. For unequal sets, choose \(z\) in their symmetric difference;
the inclusion test at \(q=p\) fails by the lower membership calculation.
These are computations of the recursion, not an appeal to generic
existence. Together with reflexivity for arbitrary names, a pair
\(\langle\check s,s\rangle\in\dot G\) makes every strengthening
of \(s\) force \(\check s\in\dot G\). This justifies both
atomic assertions in \cref{def:p06-ftr-generic-boundary}.
\end{proof}
\end{annexitem}

\begin{annexitem}{annex:p06.ftr.truth}{thm:p06-ftr-truth}{The finite formula induction}
\begin{proof}
The atomic equivalences are \cref{thm:p06-ftr-atomic-truth}; falsehood is
false in the nonempty structure and has no forcing condition. Proceed through
the expanded primitive syntax, maintaining both soundness of each forcing
condition in \(G\) and the existence of a forcing condition in \(G\)
for each true formula. For conjunction, soundness follows for each conjunct.
Conversely, two true conjuncts give two conditions in \(G\); directedness
gives one common strengthening, and monotonicity forces both there.

If \(p\in G\) forces \(\neg\theta\) and \(\theta\) were true, the
induction hypothesis would give \(q\in G\) forcing \(\theta\). A common
strengthening of \(p,q\) contradicts the negation clause. If instead
\(\neg\theta\) is true, \(G\) meets the ground decision set for
\(\theta\). Its positive alternative contradicts soundness of \(\theta\),
so its negative alternative supplies the required condition.

If \(p\in G\) forces \(\exists x\theta(x,\bar\tau)\), meet its
ground dense witness set below \(p\), using the local dense-set argument
in \cref{annex:p06.ftr.atomic-truth}. The meeting condition \(q\in G\)
has a ground name \(\sigma\) with \(q\forces\theta(\sigma,\bar\tau)\).
Soundness of the matrix gives the witness \(\sigma^G\) in the extension.
Conversely, any witness in \(M[G]\) has some ground name \(\sigma\).
Truth of its matrix gives \(q\in G\) forcing that instance. Every
strengthening of \(q\) still forces it, so its cone is contained in
the existential witness set and \(q\) forces the existential.
The set being met is a subset of \(P\) formed by ground Separation with
the fixed defining formula, not a proper-class set of witnesses.

For each of the other active constructors use its declared finite expansion.
Classical satisfaction gives
\(\varphi\lor\psi\) iff \(\neg(\neg\varphi\land\neg\psi)\),
\(\varphi\to\psi\) iff \(\neg(\varphi\land\neg\psi)\), and
\(\forall x\theta\) iff \(\neg\exists x\neg\theta\). The previous
primitive cases therefore prove both directions for each constructor.
These expansions have fresh bound variables and keep exactly the same
parameter assignments. Their forcing clauses were independently established
in \cref{thm:p06-ftr-density}. In particular the true universal direction
has not been inferred by intersecting one condition for each name; it uses
the negation/existential induction. This proves the full formula theorem
before any ZFC axiom is asserted in \(M[G]\).
\end{proof}
\end{annexitem}

\begin{annexitem}{annex:p06.ftr.zfc}{thm:p06-ftr-zfc}{Every axiom, including the two set-bounding arguments}
\begin{proof}
Write \(K=M[G]\). All statements about \(K\) below use its ordinary
set-sized membership satisfaction, already available independently of the
axioms to be verified.

\paragraph{Empty Set, Extensionality, Foundation and Infinity.}
Canonical names put \(\varnothing,\omega\) in \(K\); the latter is the
actual \(\omega\) since the transitive ZFC ground has the standard finite
ordinals and least inductive set. It remains inductive in \(K\), since
all finite successors and their memberships are unchanged. If \(a,b\in K\)
have the same members as seen in \(K\), transitivity ensures that every
actual member of either is in \(K\), so ambient Extensionality gives
\(a=b\). For nonempty \(a\in K\), ambient Foundation gives \(x\in a\)
with \(x\cap a=\varnothing\). Transitivity places \(x\) in \(K\),
and its bounded disjointness assertion is unchanged. Thus Foundation holds.

\paragraph{Pairing and Union.}
For names \(\sigma,\tau\), the ground set
\(U(\sigma,\tau)=\{\langle\sigma,1_P\rangle,
\langle\tau,1_P\rangle\}\) is a name with value
\(\{\sigma^G,\tau^G\}\). Applying \(U\) twice gives the name
\(O(\sigma,\tau)=U(U(\sigma,\sigma),U(\sigma,\tau))\) for their
Kuratowski ordered pair. For the union of \(\sigma^G\), form
\[
 L_\sigma=\{\langle\pi,q\rangle:
   \exists\langle\tau,p\rangle\in\sigma\
   \exists\langle\pi,r\rangle\in\tau\ (q\leq_P p,r)\}.
\]
The possible first coordinates belong to
\(\bigcup_{\tau\in A_\sigma}A_\tau\), a ground set by Replacement
and Union; Separation on its product with \(P\) forms \(L_\sigma\).
Every displayed entry has a lower-rank name coordinate, so it is a name.
An entry with \(q\in G\) activates both \(p,r\) and yields an element
of \(\bigcup\sigma^G\). Conversely, an element of that union has these
two activating tags in \(G\); directedness supplies their common
strengthening \(q\in G\), giving an entry of \(L_\sigma\).
Thus its value is exactly the union, not just a superset.

\paragraph{Every Separation instance.}
Fix a standard formula \(\theta(x,\bar y)\), a name \(\sigma\), and
ground names \(\bar\tau\) for its parameters. Ground Separation forms
\[
 R=\{\langle\pi,q\rangle:\exists\langle\pi,p\rangle\in\sigma\
       (q\leq_P p\land q\forces\theta(\pi,\bar\tau))\}
       \subseteq A_\sigma\times P.
\]
Definability of this fixed forcing formula makes this an instance of
Separation in \(\mathcal M\). Soundness and the activating tag show
\(R^G\subseteq\{a\in\sigma^G:K\models\theta(a,\bar\tau^G)\}\).
For the reverse inclusion choose an activating entry for \(a\), then
use the truth lemma for its matrix and take a common strengthening in
\(G\) of the two conditions. That entry belongs to \(R\).
This proves equality and handles parameters from anywhere in \(K\).

\paragraph{Collection in the ground, proved from Replacement.}
For a ground set \(I\) and a fixed ground formula \(B(i,z)\) with a
witness for each \(i\in I\), let \(\delta_i\) be the least ordinal
\(\delta\) for which some \(z\in V_\delta^M\) satisfies \(B(i,z)\).
It exists by the rank construction in
\cref{annex:p06.ftr.name-foundations}; leastness is a functional formula.
Ground Replacement collects the \(\delta_i\). Take
\(\beta=\sup_{i\in I}(\delta_i+1)\). Ground Separation on the set
\(V_\beta^M\) collects all its \(B\)-witnesses for indices in \(I\).
This set contains a witness for every index. The empty-index case uses
the empty witness set. Thus Collection at this fixed formula is derived,
not an extra axiom or a choice of a class-sized family.

\paragraph{Every Replacement instance.}
Suppose \(K\models\forall x\in\sigma^G\exists!y\theta(x,y,\bar\tau^G)\).
Inside \(\mathcal M\), use Separation to form
\[
 I=\{(\pi,q)\in A_\sigma\times P:
          \exists\nu\ (\mathsf{Name}_P(\nu)\land
                            q\forces\theta(\pi,\nu,\bar\tau))\}.
\]
Apply the ground Collection just proved to obtain a ground set \(B\) of
names such that for every \((\pi,q)\in I\), some \(\mu\in B\)
satisfies \(q\forces\theta(\pi,\mu,\bar\tau)\). Namehood can be
included in the witness formula, so every collected member is a name.
The name \(C=\{\langle\mu,1_P\rangle:\mu\in B\}\) is a ground set.
For any \(a\in\sigma^G\), choose a representing \(\pi\in A_\sigma\)
and a name \(\nu\) for its unique image. The truth lemma gives
\(q\in G\) forcing the matrix. Thus \((\pi,q)\in I\), and its
collected witness \(\mu\) has value in \(C^G\). Soundness makes that
value satisfy the same matrix; the assumed uniqueness identifies it
with the desired image. Therefore all images are bounded by the set
\(C^G\in K\). Apply the already proved Separation instance in \(K\)
with condition \(\exists x\in\sigma^G\theta(x,y,\bar\tau^G)\) to
remove any irrelevant collected values. The result is exactly the image.
No Collection or Replacement in the extension was used to build its bound.

\paragraph{Power Set, with a ground set of subset names.}
Ground Power Set supplies \(B=\mathcal P^M(A_\sigma\times P)\).
Each member is a name since its first coordinates are subnames of
\(\sigma\). The name \(C=\{\langle\nu,1_P\rangle:\nu\in B\}\)
therefore belongs to \(M\). For any \(\mu^G\in K\) with
\(\mu^G\subseteq\sigma^G\), form inside \(M\)
\[
 \nu=\{\langle\pi,q\rangle\in A_\sigma\times P:q\forces\pi\in\mu\}.
\]
It belongs to \(B\). A member of \(\nu^G\) belongs to \(\mu^G\)
by soundness. Conversely, every \(a\in\mu^G\) belongs to \(\sigma^G\)
and thus has a representative \(\pi\in A_\sigma\); the atomic truth
lemma gives \(q\in G\) forcing \(\pi\in\mu\), hence \(a\in\nu^G\).
Thus \(\nu^G=\mu^G\). All subsets of \(\sigma^G\) which belong to
\(K\) are consequently members of \(C^G\). Apply Separation in \(K\)
to \(C^G\) with condition \(y\subseteq\sigma^G\); this gives precisely
the internal power set. No ambient power set of \(K\) was assumed in \(M\).

\paragraph{Choice, after the set-existence axioms.}
Use Choice in \(\mathcal M\) to enumerate \(A_\sigma\) by an ordinal
\(\alpha\in M\), say \(e:\alpha\to A_\sigma\) is a bijection.
Ground Replacement forms
\[
 F=\{\langle O(\check\gamma,e(\gamma)),1_P\rangle:\gamma<\alpha\}.
\]
Its evaluated value is the graph of a function with domain \(\alpha\)
whose range contains \(\sigma^G\). This graph is in \(K\); ordinal
and canonical-name agreement make its first coordinates the actual ordinals.
Using the now proved set axioms, Separation on \(\sigma^G\times\alpha\)
forms the graph taking each \(a\in\sigma^G\) to the least
\(\gamma<\alpha\) with \(F^G(\gamma)=a\). It is an injection into
the ordinal and hence well-orders \(\sigma^G\). The empty set has the
empty well-order. Since every set in \(K\) is well-orderable, for a set
of nonempty sets well-order its union and use Replacement to choose each
factor's least member. This supplies its choice function in \(K\).
All axioms on the active ZFC list, and arbitrary fixed schema instances,
have now been verified. None required external countability of \(M\).
\end{proof}
\end{annexitem}

\begin{annexitem}{annex:p06.ftr.semantic}{thm:p06-ftr-semantic}{Nonvacuous semantic forcing at the countable-ground hypothesis}
\begin{proof}
When \(M\) is externally countable, its dense subsets of \(P\) have
an ambient enumeration \(D_0,D_1,\ldots\). For any \(q\in P\), choose
\(p_0=q\) and recursively \(p_{n+1}\leq_P p_n\) in \(D_n\),
using an ambient enumeration of the countable set \(P\) to take the
first suitable condition. Density guarantees a witness at each stage.
The set \(H=\{r\in P:\exists n\ p_n\leq_P r\}\) is nonempty,
upward closed and directed, contains \(q\), and meets every \(D_n\).
Thus a generic through every condition is actually constructed.

Recursive soundness gives the forward semantic implication for each such
generic. If \(p\) does not recursively force the formula, density closure
gives \(q\leq_P p\) forcing its negation. Construct a generic \(H\ni q\)
as above. It also contains \(p\), and soundness of the negation makes
the formula false in \(M[H]\). This proves the converse and excludes
vacuity. For an uncountable transitive ground with a particular generic
already supplied, the truth and ZFC-transfer theorems still apply; this
additional semantic characterization is not asserted without generics
through every relevant countercondition.
\end{proof}
\end{annexitem}
